\chapter{RQ3: Transfer Across Devices}
\label{ch:rq3}
\section{Setup}
Originally this chapter was to test a lab baseline on mock remote interviews on
students' own laptops. With classroom data only (\cref{sec:naturalistic}), the same
question is asked of the devices students actually used: many worked both on lab PCs
and on their own laptops (\cref{sec:devices}). Three comparisons, within students only:
\begin{enumerate}
  \item \emph{Drift.} The robust distance between window-feature medians of two sessions
    of the same student on the same device class, on different classes, and of two
    students on the same question.
  \item \emph{Transfer.} False alarms per active hour and splice detection when each
    session is scored against a baseline from the other device class, against the
    device-specific baseline (\texttt{--by-device}) and against a pooled one.
  \item \emph{Which features move.} Per-feature shift between device classes, to see
    whether timing (key-pair latency, hold time) or process (pauses, build order) moves.
\end{enumerate}
The pilot's first look already showed more flags on a student's less usual device
(9.9 against 5.7 per active hour, \cref{tab:pilot}).
\todo{Results on the classroom test split.}
\section{Results}
\section{Summary}
